% !TEX program = xelatex
% 问题一公式清单：无线电干扰源交会定位（定位区域直径计算与覆盖判定）
% 编译方式：xelatex 问题一公式.tex（需 ctex 宏包，随 TeX Live / MiKTeX 自带）
% 公式按论文 5.1.0 ~ 5.1.8 小节顺序编号
\documentclass[11pt]{ctexart}
\usepackage{amsmath, amssymb}
\usepackage{geometry}
\geometry{a4paper, margin=2.5cm}
\setlength{\parindent}{0pt}

\begin{document}

\begin{center}
{\LARGE\bfseries 问题一公式清单}\\[6pt]
{\normalsize 定位区域直径计算与覆盖判定（LaTeX 排版）}
\end{center}

\section*{5.1.0 数据与坐标约定}

角度差函数（避免角度跨越零点时的环绕错误）：
\begin{equation}
\mathrm{angdiff}(\varphi,\theta)
= \big((\varphi-\theta+180^\circ) \bmod 360^\circ\big) - 180^\circ
\in (-180^\circ,\; 180^\circ]
\end{equation}

单位方向向量与误差半角：
\begin{equation}
\mathbf{u}(\varphi) = (\cos\varphi,\ \sin\varphi), \qquad
\varepsilon = 1^\circ
\end{equation}

1° 误差的横向偏移（表 5-2）：
\begin{equation}
\Delta = d \cdot \tan 1^\circ
\end{equation}

参数对直径的影响（表 5-1）：
\begin{equation}
D \propto \varepsilon \ (\text{小角近似}), \qquad
D \propto \frac{1}{\sin\varphi} \ (\varphi \to 90^\circ \text{ 时最优})
\end{equation}

\section*{5.1.3 角域与定位区域的表示}

定义 5.1.1（角域）：
\begin{equation}
W(S,\theta,\varepsilon)
= \left\{ S + t\,\mathbf{u}(\varphi) : t \ge 0,\ \bigl|\mathrm{angdiff}(\varphi,\theta)\bigr| \le \varepsilon \right\}
\end{equation}

角域的边界方向：
\begin{equation}
\mathbf{u}_- = \mathbf{u}(\theta-\varepsilon), \qquad
\mathbf{u}_+ = \mathbf{u}(\theta+\varepsilon)
\end{equation}

命题 5.1（角域的两半平面形式）：
\begin{equation}
W = \left\{ X : \operatorname{cross}(\mathbf{u}_-,\, X-S) \ge 0 \right\}
\cap \left\{ X : \operatorname{cross}(\mathbf{u}_+,\, X-S) \le 0 \right\}
\end{equation}

二维叉积：
\begin{equation}
\operatorname{cross}(\mathbf{a},\mathbf{b}) = a_x b_y - a_y b_x
\end{equation}

定义 5.1.2（定位区域）：
\begin{equation}
L = \bigcap_{i=1}^{n} W_i
\end{equation}

\section*{5.1.4 定位区域的三个定理}

定理 5.2（顶点枚举与过滤）：
\begin{equation}
P \text{ 是 } L \text{ 的顶点}
\iff t \ge 0 \ \land\ s \ge 0 \ \land\ \forall k:\ P \in W_k
\end{equation}

定理 5.3（有界性判据）：
\begin{equation}
L \text{ 有界}
\iff \bigcap_{i=1}^{n} [\theta_i - \varepsilon,\ \theta_i + \varepsilon] = \varnothing
\end{equation}

定理 5.3 证明中的回收锥：
\begin{equation}
\operatorname{rec}(W) = \left\{ \mathbf{d} : W + t\mathbf{d} \subseteq W,\ \forall t \ge 0 \right\}, \qquad
\operatorname{rec}\!\Big(\bigcap_i W_i\Big) = \bigcap_i \operatorname{rec}(W_i)
\end{equation}

\section*{5.1.7 覆盖判定与 Jung 界验证}

命题 5.2（等边三角形反例，$D=a$）：
\begin{equation}
\lvert CM \rvert = \frac{\sqrt{3}}{2}\,a \approx 0.866\,a > 0.5\,a = \frac{D}{2}
\end{equation}

定理 5.4（Thales 圆内判据）：
\begin{equation}
X \text{ 在以 } PQ \text{ 为直径的圆内（含边界）} \iff \angle PXQ \ge 90^\circ
\end{equation}

定理 5.4 证明中的恒等式（$M = \tfrac{P+Q}{2}$）：
\begin{equation}
\|X-M\|^2 - \Big(\frac{D}{2}\Big)^2
= (X-P)\cdot(X-Q)
= \|P-X\|\,\|Q-X\|\cos\angle PXQ
\end{equation}

推论 5.3（可计算判据，$V(L)$ 为 $L$ 的顶点集）：
\begin{equation}
\text{以 } PQ \text{ 为直径的圆覆盖 } L
\iff \forall X \in V(L):\ \angle PXQ \ge 90^\circ
\end{equation}

定理 5.5（Jung 定理）：
\begin{equation}
R_{\mathrm{MEC}} \le \frac{D}{\sqrt{3}}
\end{equation}

定理 5.5 证明（正弦定理，$\alpha$ 为内接三角形最大角）：
\begin{equation}
D \ge s = 2R\sin\alpha \ge 2R\sin 60^\circ = \sqrt{3}\,R
\ \Longrightarrow\ R \le \frac{D}{\sqrt{3}}
\end{equation}

双侧夹逼：
\begin{equation}
\frac{D}{2} \ \le\ R_{\mathrm{MEC}}\ \le\ \frac{D}{\sqrt{3}} \approx 0.5774\,D
\end{equation}

覆盖缺口比：
\begin{equation}
\gamma = \frac{R_{\mathrm{MEC}}}{D/2} \in \left[1,\ \frac{2}{\sqrt{3}} \approx 1.1547\right]
\end{equation}

\section*{5.1.8 与问题二的联系}

Fisher 信息矩阵：
\begin{equation}
M = J^{\mathsf{T}}J
\end{equation}

两站情形的行列式：
\begin{equation}
\det M = \frac{\sin^2\varphi}{\sigma^4 d_1^2 d_2^2}
\end{equation}

均方根位置误差：
\begin{equation}
E = \varepsilon\,\sqrt{\operatorname{tr}(M^{-1})}
\end{equation}

\end{document}
